\chapter{Extension multi-classes un-contre-tous}
\label{ch:ovr}

\begin{lecture}
OvR transforme le problème à quatre maisons en quatre questions binaires :
« Gryffindor contre les autres », puis Hufflepuff, Ravenclaw et Slytherin. Pour
un nouvel élève, les quatre modèles produisent quatre scores ; la maison associée
au score maximal devient la prédiction.
\end{lecture}

Pour chaque classe $k\in\{1,\ldots,K\}$, on construit la cible binaire
\[
  y_{ik}=\ind[c_i=k].
\]
Cette réduction est la stratégie one-vs-rest : un classifieur binaire par classe,
la classe considérée étant positive et toutes les autres négatives
\cite[sec.~4.1]{rifkin2004}. On ajuste ensuite indépendamment un paramètre
$\thv_k$ pour chacun de ces problèmes. On obtient les scores
$z_k(\x)=\thv_k^\top\xt$ et les sorties binaires
$q_k(\x)=\sigma(z_k(\x))$. La règle de décision est
\[
  \widehat c(\x)=\Argmax_{k\in\{1,\ldots,K\}}q_k(\x)
  =\Argmax_{k\in\{1,\ldots,K\}}z_k(\x),
  \tag{3}\label{eq:ovr-decision}
\]
car $\sigma$ est strictement croissante et commune aux $K$ modèles.
La décision par maximum des réponses binaires est la règle OvR usuelle
\cite{rifkin2004}.

L'apprentissage produit ainsi $K$ vecteurs de paramètres
$\thv_1,\ldots,\thv_K$. Pour quatre maisons, une même matrice $\X$ est associée
successivement à quatre vecteurs cibles $\y_1,\ldots,\y_4$. À la prédiction, une
observation préparée est multipliée par chacun des quatre vecteurs de paramètres ;
les quatre scores obtenus sont comparés, et l'indice du maximum fournit la maison.

\begin{table}[htbp]
\centering
\begin{tabular}{lrrr}
\toprule
classe $k$ & score $z_k$ & sortie $q_k=\sigma(z_k)$ & rang \\
\midrule
Gryffindor & $-1{,}99$ & $0{,}12$ & $3$ \\
Hufflepuff & $-2{,}31$ & $0{,}09$ & $4$ \\
Ravenclaw  & $ 0{,}90$ & $0{,}71$ & $1$ \\
Slytherin  & $-1{,}21$ & $0{,}23$ & $2$ \\
\bottomrule
\end{tabular}
\caption{Décision OvR sur une observation. Ravenclaw maximise à la fois $z_k$ et
$q_k$, la même fonction strictement croissante $\sigma$ étant appliquée aux
quatre scores. La somme des sorties vaut ici $1{,}15$.}
\label{tab:ovr-example}
\end{table}

\begin{figure}[htbp]
\centering
\resizebox{\textwidth}{!}{%
\begin{tikzpicture}[x=.105cm,y=1cm]
  \node[anchor=east] at (0,3.5){G};
  \node[anchor=east] at (0,2.5){H};
  \node[anchor=east] at (0,1.5){R};
  \node[anchor=east] at (0,.5){S};
  \foreach \yy/\v in {3.5/12,2.5/9,1.5/71,.5/23}{
    \fill[accent!70] (0,\yy-.22) rectangle (\v,\yy+.22);
    \node[anchor=west] at (\v+1,\yy){\v\,\%};
  }
  \node[font=\bfseries] at (38,4.35){sorties OvR indépendantes : somme $115\,\%$};
  \begin{scope}[xshift=9.4cm]
    \node[anchor=east] at (0,3.5){G};
    \node[anchor=east] at (0,2.5){H};
    \node[anchor=east] at (0,1.5){R};
    \node[anchor=east] at (0,.5){S};
    \foreach \yy/\v in {3.5/5,2.5/3,1.5/82,.5/10}{
      \fill[warning!70] (0,\yy-.22) rectangle (\v,\yy+.22);
      \node[anchor=west] at (\v+1,\yy){\v\,\%};
    }
    \node[font=\bfseries] at (40,4.35){softmax des scores : somme $100\,\%$};
  \end{scope}
\end{tikzpicture}}
\caption{Comparaison sur les scores de la table~\ref{tab:ovr-example}. OvR
applique une sigmoïde séparée à chaque score ; softmax normalise conjointement
les exponentielles des scores. Les deux méthodes choisissent ici Ravenclaw, mais
leurs valeurs n'ont pas la même signification probabiliste.}
\label{fig:ovr-softmax}
\end{figure}

\begin{figure}[htbp]
  \centering
  \resizebox{\textwidth}{!}{%
  \begin{tikzpicture}[
    model/.style={draw=accent,fill=accentlight,rounded corners,minimum width=39mm,
      minimum height=8mm,align=center},
    score/.style={draw=ink!55,fill=white,rounded corners,minimum width=18mm,
      minimum height=8mm,align=center},arr/.style={-{Latex},draw=ink}]
    \node[model] (x) at (0,0) {même observation $\xt$};
    \node[model] (mG) at (5.2, 2.1) {Gryffindor contre le reste};
    \node[model] (mH) at (5.2,  .7) {Hufflepuff contre le reste};
    \node[model] (mR) at (5.2, -.7) {Ravenclaw contre le reste};
    \node[model] (mS) at (5.2,-2.1) {Slytherin contre le reste};
    \node[score] (sG) at (9.0, 2.1) {$z_G=-1{,}99$};
    \node[score] (sH) at (9.0,  .7) {$z_H=-2{,}31$};
    \node[score] (sR) at (9.0, -.7) {$z_R= 0{,}90$};
    \node[score] (sS) at (9.0,-2.1) {$z_S=-1{,}21$};
    \node[model,fill=warning!18,draw=warning] (max) at (12.6,0)
      {$\Argmax_k z_k=R$\\prédiction : Ravenclaw};
    \foreach \m/\s in {mG/sG,mH/sH,mR/sR,mS/sS}{
      \draw[arr] (x.east)--(\m.west);
      \draw[arr] (\m.east)--(\s.west);
      \draw[arr] (\s.east)--(max.west);
    }
  \end{tikzpicture}}
  \caption{Architecture OvR complète pour les quatre maisons. La même
  observation traverse les quatre modèles binaires, qui produisent quatre scores,
  ensuite comparés.}
\label{fig:ovr-architecture}
\end{figure}

\begin{procedure}{Ajuster et utiliser un modèle OvR}
\textbf{Entrées d'apprentissage :} $\X\in\R^{m\times(n+1)}$ et les classes
$c_i$. \textbf{État produit :} $K$ vecteurs $\thv_k$ associés à $K$ libellés.
\textbf{Sortie de prédiction :} un libellé de classe.
\begin{enumerate}
  \item pour la classe $k$, construire $y_{ik}=1$ si $c_i=k$, sinon $0$ ;
  \item ajuster $\thv_k$ avec $\X$ et ce vecteur cible, puis répéter pour les $K$ classes ;
  \item stocker les paramètres dans un ordre explicite, par exemple
        G, H, R, S, et conserver cet ordre avec le modèle ;
  \item pour une nouvelle ligne $\xt$, calculer
        $(z_G,z_H,z_R,z_S)=(\thv_G^\top\xt,\ldots,\thv_S^\top\xt)$ ;
  \item prendre la position de la plus grande valeur ; par exemple
        $(-1{,}99,-2{,}31,0{,}90,-1{,}21)$ renvoie la troisième position, donc Ravenclaw ;
  \item vérifier que quatre scores ont été produits et qu'aucun n'est non fini.
\end{enumerate}
\end{procedure}

\begin{vigilance}[Scores, probabilités binaires et probabilités multiclasses]
$q_k$ estime la probabilité de la cible binaire « classe $k$ contre le reste »
dans le modèle $k$. Les $K$ modèles étant ajustés séparément, chaque $q_k$ décrit
sa propre expérience binaire et la somme $\sum_k q_k$ prend une valeur libre. Ces
nombres servent au classement de l'équation~\eqref{eq:ovr-decision} ; leur usage
comme vecteur de probabilités multiclasses \rsym{calibration}{calibrées} demande une recalibration
préalable \cite{zadrozny2002}. Une régression logistique multinomiale à fonction
softmax fournit directement une \rsym{categorielle}{loi catégorielle} lorsque cet objectif
probabiliste est central \cite{hastie2009}.
\end{vigilance}

\begin{resultat}[Règle de décision OvR]
La décision OvR tient en deux opérations : calculer les quatre scores $z_k$, puis
retenir $\Argmax_k z_k$. La section suivante présente softmax comme un modèle
probabiliste multiclasses distinct, avec sa propre estimation.
\end{resultat}

La fonction \rsym{softmax}{softmax} mentionnée dans la comparaison est
\[
  \Pp(C=k\mid\x)
  =\frac{e^{z_k}}{\sum_{r=1}^K e^{z_r}}.
\]
Son dénominateur $\sum_r e^{z_r}$ couple les $K$ classes et impose une somme
égale à $1$. Dans OvR, le dénominateur $1+e^{-z_k}$ dépend du seul score $z_k$.
Cette différence algébrique fonde la différence d'interprétation ; sur
l'exactitude, les deux méthodes se départagent au cas par cas.

Les négatifs d'un problème OvR réunissent les trois autres maisons : ils sont
plus nombreux et plus hétérogènes que les positifs. Chaque sous-modèle affiche
donc une exactitude binaire élevée par construction. Le jugement porte sur la
décision multiclasses finale et sur les performances par maison.

\begin{figure}[htbp]
\centering
\begin{tikzpicture}[scale=1.15]
  \fill[accent!18] (-4,-3) rectangle (0,3);
  \fill[warning!18] (0,-3) rectangle (4,0);
  \fill[blue!12] (0,0) rectangle (4,3);
  \fill[orange!14] (-4,-3) rectangle (0,0);
  \draw[-{Latex},thick] (-4.2,0)--(4.4,0) node[right]{$x_1$};
  \draw[-{Latex},thick] (0,-3.2)--(0,3.4) node[above]{$x_2$};
  \draw[ink,very thick] (-4,0)--(4,0);
  \draw[ink,very thick] (0,-3)--(0,3);
  \node[align=center] at (-2,1.5){Gryffindor\\$z_G$ maximal};
  \node[align=center] at (1.6,1.5){Ravenclaw\\$z_R$ maximal};
  \node[align=center] at (-2,-1.5){Hufflepuff\\$z_H$ maximal};
  \node[align=center] at (2,-1.5){Slytherin\\$z_S$ maximal};
  \fill[ink] (3.0,.65) circle (2.2pt);
  \node[draw=ink!55,fill=white,rounded corners=1pt,
        inner sep=3pt,align=center] (obs) at (3.15,2.45)
        {observation $\x$\\à classer};
  \draw[-{Latex},ink!70] (obs.south)--(3.02,.75);
\end{tikzpicture}
\caption{Régions de décision OvR dans un espace à deux caractéristiques. Chaque
point rejoint la région dont le score domine. Les frontières entre deux régions
vérifient $z_k(\x)=z_r(\x)$ et restent linéaires tant que les scores le sont.}
\label{fig:ovr-regions}
\end{figure}

La règle du maximum s'applique dans deux situations remarquables. Quand les
quatre scores sont négatifs, elle retient le moins négatif. Quand plusieurs
scores sont fortement positifs, elle retient le plus grand. Dans les deux cas
elle produit une maison. Une politique de rejet forme une décision
supplémentaire, définie et validée à part.

Exemple de rejet : laisser la maison vide si le meilleur score descend sous $-2$,
ou si l'écart entre les deux meilleurs scores reste sous $0{,}1$. Ces seuils
viennent du protocole et introduisent leur propre arbitrage : davantage de rejets
réduit les mauvaises attributions et laisse davantage d'élèves sans réponse. Ils
se choisissent sur validation d'après le coût respectif d'un rejet et d'une
erreur, puis s'évaluent sur test.

\begin{resultat}[État conservé par le modèle OvR]
Avec $K$ classes, l'apprentissage conserve $K$ vecteurs de paramètres et l'ordre
qui associe chaque vecteur à son libellé. Une prédiction produit un vecteur de
$K$ scores dans ce même ordre, dont l'indice du maximum se convertit en libellé.
Cette correspondance appartient au modèle entraîné.
\end{resultat}
